% #############################################################################
% This is Appendix C
% !TEX root = ../main.tex
% #############################################################################
\chapter[Isolated Reward-Based Continuation Ablations]{Isolated Reward-Based\\Continuation Ablations}
\label{chap:appendix-rl-ablation}

This appendix reports isolated reward-design diagnostics on the fully fine-tuned
270M-parameter T5Gemma 2 model. They are separate from the Stage C
continuations with the 4B-parameter model described in
\Cref{sec:method-stage-c,sec:rewriting-results}. Four independent diagnostics examine
the base normalized BLEU/TER reward and three modifications involving classifier
feedback, token conservatism, and candidate-group construction. Both
encoder-controlled and decoder-controlled formulations are examined.

\section{Experimental Setup}

Each diagnostic begins independently from the selected supervised Stage B
checkpoint of the 270M-parameter model for its control formulation. These
checkpoints follow Stage A adaptation on OpenSubtitles and Stage B
specialization on GPT-Wikipedia under the supervised protocol selected in
\Cref{sec:smaller-model-protocol-study}. No reward variant is continued from
another.

The two control formulations start from different supervised checkpoints, so
each diagnostic is interpreted relative to its own Stage B baseline rather than
compared in absolute terms.

All continuations use the same 1200-step optimization budget. Evaluation follows
the 270M-parameter protocol-selection comparison, reporting macro-F1 for
identification and aggregate corpus BLEU and TER for rewriting on FRMT and the
Golden Collection.

\section{Reward Variants}

The base reward uses the normalized BLEU and TER components defined in
\Cref{sec:method-stage-c}. For each generated candidate, sentence-level BLEU
and TER are computed against the reference rewrite and normalized relative to
the scores obtained by leaving the source unchanged. The base reward weights both components equally and penalizes source copies by 0.10 when the reference requires a change.

In the decoder-controlled diagnostic, the generated \texttt{PT}/\texttt{BR}
control label is also checked before the sentence-level reward is used. An incorrect first label sets the candidate reward to zero, whereas main
decoder-controlled Stage C uses a weighted first-label component.

The classifier-reward variant adds an auxiliary variety-identification score.
Each generated rewrite is passed back through the model using the identification
formulation for the same control setup, and the probability assigned to the
intended variety of the rewritten sentence provides the classifier score. This
measures whether the generated sentence itself exhibits the requested variety.
Under decoder control, this component is separate from the binary first-label
gate, which checks only the generated control token.

The token-conservatism variant adds a token-level score. After lowercasing and
tokenizing the source, candidate, and reference, each generated token receives one of three scores.
A token absent from the reference receives 0. A token present in both the
reference and source receives 0.5. A token present in the reference but absent
from the source receives 1.0. The sequence score averages these token scores.

For both auxiliary-reward variants, normalized BLEU and TER each receive a
weight of 0.45, while the added component receives a weight of 0.10.

The TER-diverse grouping variant changes candidate-group construction rather
than the reward formula. For each input, a pool of 12 candidates is first
sampled and their TER against the reference is computed. After duplicate
normalized candidate texts are removed, the remaining candidates are ordered
by TER and up to four are selected at evenly spaced quantile positions across that
ordering, producing a group that covers different levels of distance from the
reference. The reference candidate is then appended, so the group-relative
optimization normally compares four selected samples with the reference
candidate. The selected candidates use the same normalized BLEU/TER reward as
the base configuration. This diagnostic therefore changes the examples compared
within the group-relative update rather than adding another reward component.

\section{Results}

\Cref{tab:appendix-rl-d,tab:appendix-rl-e} report the decoder-controlled and
encoder-controlled diagnostics, respectively.

\begin{table}[htbp]
\centering
\thesistablefont
\setlength{\tabcolsep}{4pt}
\caption{Decoder-controlled 270M isolated Stage C ablations.}
\label{tab:appendix-rl-d}
\begin{tabularx}{\textwidth}{@{}Ycccccc@{}}
\toprule
\textbf{Configuration} &
\multicolumn{3}{c}{\textbf{FRMT}} &
\multicolumn{3}{c}{\textbf{Golden Collection}} \\
\cmidrule(lr){2-4} \cmidrule(lr){5-7}
& \shortstack{\textbf{Macro-F1}\\\textbf{(\%)}} & \shortstack{\textbf{BLEU}\\\textbf{(0--100)}} & \textbf{TER}
& \shortstack{\textbf{Macro-F1}\\\textbf{(\%)}} & \shortstack{\textbf{BLEU}\\\textbf{(0--100)}} & \textbf{TER} \\
\midrule
Supervised Stage B start & \textbf{70.76} & \textbf{39.13} & \textbf{0.4788} & \textbf{67.63} & 86.57 & 0.0980 \\
\midrule
Base reward & 68.76 & 39.11 & 0.4800 & 67.15 & 86.88 & 0.0962 \\
Classifier reward & 69.54 & 39.04 & 0.4806 & 67.57 & 86.96 & 0.0969 \\
Token conservatism & 69.49 & 39.07 & 0.4804 & 67.43 & 86.58 & 0.1007 \\
TER-diverse grouping & 68.24 & 39.12 & 0.4799 & 66.70 & \textbf{87.19} & \textbf{0.0949} \\
\bottomrule
\end{tabularx}
\end{table}

\begin{table}[htbp]
\centering
\thesistablefont
\setlength{\tabcolsep}{4pt}
\caption{Encoder-controlled 270M isolated Stage C ablations.}
\label{tab:appendix-rl-e}
\begin{tabularx}{\textwidth}{@{}Ycccccc@{}}
\toprule
\textbf{Configuration} &
\multicolumn{3}{c}{\textbf{FRMT}} &
\multicolumn{3}{c}{\textbf{Golden Collection}} \\
\cmidrule(lr){2-4} \cmidrule(lr){5-7}
& \shortstack{\textbf{Macro-F1}\\\textbf{(\%)}} & \shortstack{\textbf{BLEU}\\\textbf{(0--100)}} & \textbf{TER}
& \shortstack{\textbf{Macro-F1}\\\textbf{(\%)}} & \shortstack{\textbf{BLEU}\\\textbf{(0--100)}} & \textbf{TER} \\
\midrule
Supervised Stage B start & 71.73 & \textbf{27.06} & \textbf{0.6204} & 65.43 & \textbf{65.53} & \textbf{0.2996} \\
\midrule
Base reward & 71.90 & 13.62 & 0.7525 & 65.77 & 38.21 & 0.5154 \\
Classifier reward & 71.81 & 15.67 & 0.7319 & 66.29 & 44.25 & 0.4704 \\
Token conservatism & \textbf{72.28} & 16.27 & 0.7268 & \textbf{66.34} & 42.79 & 0.4808 \\
TER-diverse grouping & 72.22 & 13.18 & 0.7571 & 66.14 & 38.96 & 0.5118 \\
\bottomrule
\end{tabularx}
\end{table}

\looseness=-1
For the decoder-controlled diagnostics, FRMT rewriting remains close to the
supervised Stage B starting point across all four variants, with a largest
absolute BLEU change of only 0.09 points. However, on the Golden Collection,
TER-diverse grouping gives the best rewriting result among the diagnostics,
reaching 87.19 BLEU and 0.0949 TER compared with 86.57 BLEU and 0.0980 TER at
the supervised starting point. Its macro-F1 is lower, at 66.70 compared with
67.63. Token conservatism does not improve Golden Collection TER, reaching
0.1007 compared with 0.0980 at the starting point.

\looseness=-1
The encoder-controlled diagnostics show a different pattern. Every
reward-based continuation substantially reduces rewriting quality relative to
the supervised Stage B starting point. On FRMT, the best diagnostic BLEU is
16.27 under token conservatism, compared with 27.06 at the starting point. On
the Golden Collection, classifier reward gives the highest diagnostic BLEU at
44.25, compared with 65.53 before continuation. Identification changes are much
smaller than the rewriting changes, with token conservatism giving the highest
macro-F1 among the diagnostics on both FRMT and the Golden Collection.

\looseness=-1
Taken together, no tested reward modification improves rewriting consistently
across both control formulations. Decoder-controlled continuation remains close
to its supervised starting point and yields modest Golden Collection gains
under some variants, with TER-diverse grouping giving the best
decoder-controlled rewriting result. In contrast, every encoder-controlled
continuation substantially reduces rewriting despite much smaller identification
changes. The classifier and token-conservatism components provide no consistent
advantage across formulations, so no tested modification is generally
preferable.
